% SPDX-License-Identifier: CC-BY-SA-4.0
% Author: Matthieu Perrin
% Part: <Nom de la partie>
% Section: <Nom de la section>
% Sub-section: <Nom de la sous-section>  % (facultatif, laisser vide si non utilisé)
% Frame: <Titre de la slide>

\begingroup

\SetKwFunction{Decide}{decide}

\begin{frame}{Retour sur l'Entscheidungsproblem}

  Soit $\langle \Phi, \Pi, \valid \rangle$ le système de preuve formel des mathématiques usuelles

  \Probleme{Entscheidungsproblem}{
    Une proposition $\varphi \in \Phi$ 
  }{
    Est-ce que $\varphi$ est démontrable ($\exists \pi \in \Pi,   \pi \valid \varphi$) ?
  }
  
  \begin{block}{Théorème -- Turing, Church (1936)}
    Le \textsc{Entscheidungsproblem} est indécidable. 
  \end{block}

  \pause 
  \begin{block}{Démonstration}
    Si le \structure{\textsc{Entscheidungsproblem}} était décidable, \structure{$\textsc{Universal}_{\langle \mathcal{L}, \llbracket \cdot \rrbracket \rangle}$} aussi :

    \begin{algorithm}[H]
      \Fun{$\Decide_{\textsc{Universal}}(P \in L, u \in \Omega^\star) \in \mathbb{B}$}{
        \Return $\Decide_{\textsc{Entscheidungsproblem}}(``u \in \llbracket P \rrbracket")$;
      }
    \end{algorithm}

    Or \structure{$\textsc{Universal}_{\langle \mathcal{L}, \llbracket \cdot \rrbracket \rangle}$} est indécidable, donc \structure{\textsc{Entscheidungsproblem}} aussi. 
  \end{block}
  
\end{frame}

\endgroup
